\documentclass[11pt]{article}
\usepackage[margin=25mm]{geometry}
\usepackage{fontspec}
\setmainfont{Latin Modern Roman}
\usepackage{amsmath,amssymb,mathtools}
\usepackage{unicode-math}
\setmathfont{Latin Modern Math}
\usepackage{xurl,hyperref}
\hypersetup{colorlinks=true,urlcolor=blue}
\setcounter{secnumdepth}{0}
\setlength{\parindent}{0pt}
\setlength{\parskip}{5pt}
\setlength{\emergencystretch}{3em}
\newcommand{\sref}[2]{\hyperlink{source:#1:#2}{[S#2]}}
\newcommand{\cataloguescope}[1]{\paragraph{Recorded target.}#1}
\begin{document}
\section{MLC Conjecture}\label{statement}

\subsection{Definitions and mathematical statement}

\paragraph{Shared notation.}$\mathbb N=\{1,2,\ldots\}$, and $\mathbb Z,\mathbb Q,\mathbb R,\mathbb C$ denote the integers, rationals, reals and complex numbers. Any other conventions are specified locally.


For $c\in\mathbb C$, define $z_0(c)=0$ and $z_{m+1}(c)=z_m(c)^2+c$ for integers $m\ge0$. The Mandelbrot set is
\[
 \mathcal M=\{c\in\mathbb C:\sup_{m\ge0}|z_m(c)|<\infty\},
\]
with the topology inherited from $\mathbb C$. A neighborhood below means a relatively open set in $\mathcal M$.
\paragraph{Statement recorded in the supplied source.}
Is $\mathcal M$ locally connected? Explicitly, for every $c\in\mathcal M$ and every relatively open neighborhood $U$ of $c$ in $\mathcal M$, does there exist a connected relatively open set $V$ such that $c\in V\subseteq U$? \sref{statement}{2}
\subsection{Short English statement}
Does every point of the Mandelbrot set have arbitrarily small connected neighborhoods within the Mandelbrot set?
\subsection{Sources}
\begin{description}
\item[\hypertarget{source:statement:1}{S1}] ULAM.AI and the UnsolvedMath Contributors, \emph{UnsolvedMath: A Curated Collection of Open Mathematics Problems}, catalogue entry. CC BY 4.0. \url{https://huggingface.co/datasets/ulamai/UnsolvedMath}.
\item[\hypertarget{source:statement:2}{S2}] Dierk Schleicher, \emph{On Fibers and Local Connectivity of Mandelbrot and Multibrot Sets} (1999), arXiv:math/9902155v1, Sections 1--2. \url{https://arxiv.org/abs/math/9902155v1}.
\end{description}
\label{end:statement}
\end{document}
